\section{Related Work}
\label{sec:related_work}

\paragraph{CLIP-based zero-shot anomaly detection.}
WinCLIP~\cite{jeong2023winclip} compares image and window features with
normal/anomalous text prompts, and its few-shot version uses normal reference
images for visual association. PromptAD~\cite{li2024promptad},
AnomalyCLIP~\cite{zhou2024anomalyclip}, and AdaCLIP~\cite{cao2024adaclip}
improve the text side through manual, object-agnostic, or hybrid learnable
prompts. VCP-CLIP~\cite{qu2024vcpclip} injects visual context into text
embeddings, while AA-CLIP~\cite{ma2025aaclip} trains anomaly-aware anchors and
residual adapters. These methods primarily ask how to obtain better
normal/abnormal representations. Our work asks a complementary question: once
such prototypes are available, which local responses in a single test image
should be trusted as evidence for calibrating them?

\paragraph{Frequency and local-structure cues.}
Frequency cues are a relevant but already established direction. FE-CLIP~
\cite{gong2025feclip} introduces frequency-enhanced adapters for CLIP.
WMoE-CLIP~\cite{chen2026wmoeclip} decomposes CLIP features into Haar
low-frequency and high-frequency components for wavelet-enhanced prompt
learning. HarmoniAD~\cite{zhang2026harmoniad} separates local structural
signals and global semantic signals in the frequency domain. These works make
it inappropriate to claim that this paper discovers the usefulness of
frequency. The distinction is usage: \method{} does not train frequency
features into CLIP and does not treat a wavelet response as the anomaly map.
The response is used only as a boundary-aware reliability signal for patch
evidence selection.

\paragraph{Test-time adaptation and reference estimation.}
Test-time prompt tuning~\cite{shu2022tpt} and conditional prompt learning~
\cite{zhou2022cocoop} show that image-conditioned prompting can improve
vision-language generalization. In anomaly detection, WinCLIP+ uses external
normal reference images, and VCP-CLIP uses trained visual context modules. Our
setting is different from both: no target-category normal bank is accessed, and
no prompt or adapter parameter is optimized at test time. The estimated
reference is constructed from the current test image itself. This makes the
method closer to a conservative per-image scoring calibration than to prompt
learning or self-training.
